\documentclass[12pt,a4paper]{article}
\usepackage[UTF8]{ctex}
\usepackage{amsmath,amssymb,amsthm}
\usepackage{geometry}
\usepackage{bm}
\usepackage{hyperref}
\usepackage{tikz}
\usetikzlibrary{arrows,arrows.meta,positioning,calc,shapes.geometric,decorations.pathreplacing}

\geometry{margin=2.5cm}

\title{12.1.6难点详解：CoR的深入理解}
\author{第12章抓取与操作补充说明}
\date{}

\newtheorem{theorem}{定理}
\newtheorem{definition}{定义}
\newtheorem{example}{例子}

\begin{document}

\maketitle

\section{概述}

本文档专门解释12.1.6章节中的难点，包括：
\begin{enumerate}
\item 正线性组合对应的CoR线段/射线
\item 不同符号角速度的情况
\item 多个CoR标记的约束区域
\end{enumerate}

\section{基础回顾：从Twist到CoR}

\subsection{平面Twist的表示}

平面twist只有3个分量：
\begin{equation}
V = (\omega_z, v_x, v_y)
\end{equation}

\subsection{CoR的计算}

如果$\omega_z \neq 0$，CoR的位置为：
\begin{equation}
\text{CoR} = \left(-\frac{v_y}{\omega_z}, \frac{v_x}{\omega_z}\right)
\end{equation}

\textbf{关键理解}：CoR是平面上的一个\textbf{点}，表示物体绕这个点旋转。

\subsection{标记系统}

\begin{itemize}
\item \textbf{'+'}：$\omega_z > 0$（逆时针旋转）
\item \textbf{'-'}：$\omega_z < 0$（顺时针旋转）
\item \textbf{'$\pm$'}：在接触法线上（可以是滚动或滑动）
\end{itemize}

\section{难点1：两个Twist的线性组合}

\subsection{问题设置}

给定两个twist：
\begin{align}
V_1 &= (\omega_{1z}, v_{1x}, v_{1y}) \\
V_2 &= (\omega_{2z}, v_{2x}, v_{2y})
\end{align}

它们的线性组合：
\begin{equation}
V = k_1 V_1 + k_2 V_2 = (k_1 \omega_{1z} + k_2 \omega_{2z}, k_1 v_{1x} + k_2 v_{2x}, k_1 v_{1y} + k_2 v_{2y})
\end{equation}

其中$k_1, k_2 \in \mathbb{R}$可以是任意实数。

\subsection{CoR的线性组合公式}

如果$\omega_z = k_1 \omega_{1z} + k_2 \omega_{2z} \neq 0$，则：
\begin{align}
\text{CoR}(V) &= \left(-\frac{v_y}{\omega_z}, \frac{v_x}{\omega_z}\right) \\
&= \left(-\frac{k_1 v_{1y} + k_2 v_{2y}}{k_1 \omega_{1z} + k_2 \omega_{2z}}, \frac{k_1 v_{1x} + k_2 v_{2x}}{k_1 \omega_{1z} + k_2 \omega_{2z}}\right)
\end{align}

\textbf{关键性质}：当$k_1$和$k_2$变化时，$\text{CoR}(V)$在通过$\text{CoR}(V_1)$和$\text{CoR}(V_2)$的\textbf{直线}上移动。

\subsection{具体数值例子}

\textbf{例子1}：两个'+' CoR

设：
\begin{align}
V_1 &= (1, 2, 0) \quad \Rightarrow \quad \text{CoR}_1 = (0, 2) \\
V_2 &= (1, 0, 2) \quad \Rightarrow \quad \text{CoR}_2 = (-2, 0)
\end{align}

线性组合：$V = k_1 V_1 + k_2 V_2 = (k_1 + k_2, 2k_1, 2k_2)$

当$k_1 + k_2 \neq 0$时：
\begin{equation}
\text{CoR}(V) = \left(-\frac{2k_2}{k_1 + k_2}, \frac{2k_1}{k_1 + k_2}\right)
\end{equation}

当$k_1$和$k_2$变化时，CoR在通过$(0, 2)$和$(-2, 0)$的直线上移动。

\section{难点2：正线性组合（$k_1, k_2 \geq 0$）}

\subsection{为什么重要？}

正线性组合对应\textbf{twist锥}，这在接触约束中非常重要，因为约束通常形成锥形。

\subsection{情况1：相同符号的角速度}

\textbf{条件}：$\omega_{1z}$和$\omega_{2z}$同号（都是正或都是负）

\textbf{结论}：正线性组合的CoR都在两个CoR之间的\textbf{线段}上，且符号相同。

\textbf{数学证明}：

设$\omega_{1z} > 0$，$\omega_{2z} > 0$（都是'+'）

对于$k_1, k_2 \geq 0$：
\begin{align}
\omega_z &= k_1 \omega_{1z} + k_2 \omega_{2z} > 0 \quad \text{（总是正的）}
\end{align}

CoR位置可以写成：
\begin{equation}
\text{CoR}(V) = \frac{k_1}{k_1 + k_2} \text{CoR}_1 + \frac{k_2}{k_1 + k_2} \text{CoR}_2
\end{equation}

这是$\text{CoR}_1$和$\text{CoR}_2$的\textbf{凸组合}（系数非负且和为1），所以CoR在连接两个CoR的线段上。

\textbf{图示}：

\begin{center}
\begin{tikzpicture}[scale=2]
% 两个'+' CoR
\filldraw[red] (0,2) circle (2pt) node[above left] {CoR$_1$ (+)};
\filldraw[red] (-2,0) circle (2pt) node[below left] {CoR$_2$ (+)};

% 连接线段
\draw[thick,blue,line width=3pt] (0,2) -- (-2,0);
\node[blue] at (-1,1) {可行CoR区域};

% 标记几个点
\filldraw[green] (-0.5,1.5) circle (1.5pt);
\filldraw[green] (-1,1) circle (1.5pt);
\filldraw[green] (-1.5,0.5) circle (1.5pt);

\node[below] at (-1,-0.3) {所有点都是'+'标记};
\end{tikzpicture}
\end{center}

\subsection{情况2：不同符号的角速度（重点难点）}

\textbf{条件}：$\omega_{1z} > 0$（'+'），$\omega_{2z} < 0$（'-'）

\textbf{关键观察}：当$k_1$和$k_2$变化时，$\omega_z = k_1 \omega_{1z} + k_2 \omega_{2z}$的符号会改变！

\textbf{详细分析}：

\begin{enumerate}
\item \textbf{当$k_1 \omega_{1z} + k_2 \omega_{2z} > 0$时}：
  \begin{itemize}
  \item $\omega_z > 0$，CoR标记为'+'
  \item CoR在通过两个CoR的直线上
  \end{itemize}

\item \textbf{当$k_1 \omega_{1z} + k_2 \omega_{2z} < 0$时}：
  \begin{itemize}
  \item $\omega_z < 0$，CoR标记为'-'
  \item CoR也在同一条直线上
  \end{itemize}

\item \textbf{当$k_1 \omega_{1z} + k_2 \omega_{2z} = 0$时}：
  \begin{itemize}
  \item $\omega_z = 0$，纯平移
  \item CoR在无穷远处
  \end{itemize}
\end{enumerate}

\textbf{关键问题}：哪些CoR是可行的？

对于正线性组合（$k_1, k_2 \geq 0$），我们需要找到满足$\omega_z = k_1 \omega_{1z} + k_2 \omega_{2z} \geq 0$或$\leq 0$的区域。

\textbf{数学分析}：

设$\omega_{1z} = 1$，$\omega_{2z} = -1$（简化，不失一般性）

则：$\omega_z = k_1 - k_2$

\begin{itemize}
\item 当$k_1 > k_2$时：$\omega_z > 0$（'+'）
\item 当$k_1 < k_2$时：$\omega_z < 0$（'-'）
\item 当$k_1 = k_2$时：$\omega_z = 0$（平移）
\end{itemize}

\textbf{CoR位置分析}：

当$k_1$和$k_2$变化时，CoR在通过$\text{CoR}_1$和$\text{CoR}_2$的整条直线上移动。

但是：
\begin{itemize}
\item \textbf{靠近CoR$_1$的区域}：$k_1$大，$k_2$小 → $\omega_z > 0$ → '+'标记
\item \textbf{靠近CoR$_2$的区域}：$k_1$小，$k_2$大 → $\omega_z < 0$ → '-'标记
\item \textbf{中间区域}：$k_1 \approx k_2$ → $\omega_z \approx 0$ → 接近平移
\end{itemize}

\textbf{关键结论}：

正线性组合的CoR由以下部分组成：
\begin{enumerate}
\item 从CoR$_1$出发的'+' CoR射线（$k_1 > k_2$）
\item 从CoR$_2$出发的'-' CoR射线（$k_1 < k_2$）
\item 两个CoR之间的线段\textbf{不可行}（因为$\omega_z$会变号，但正组合要求$k_1, k_2 \geq 0$，所以中间段实际上对应$\omega_z$接近0的情况，即平移）
\item 无穷远处的'0'点（对应$k_1 = k_2$，纯平移）
\end{enumerate}

\textbf{图示}：

\begin{center}
\begin{tikzpicture}[scale=2]
% 两个不同符号的CoR
\filldraw[red] (0,2) circle (2pt) node[above left] {CoR$_1$ (+)};
\filldraw[blue] (-2,0) circle (2pt) node[below left] {CoR$_2$ (-)};

% 整条直线
\draw[thick,green] (-3,-1) -- (1,3);

% 标记中间不可行区域（实际上这部分对应平移）
\draw[thick,red,dashed] (-0.5,1.5) -- (-1.5,0.5);
\node[red] at (-1,1) {平移区域};

% 标记可行区域
\draw[thick,blue,->,line width=2pt] (0,2) -- (0.8,2.8) node[right] {可行 (+)};
\draw[thick,green,->,line width=2pt] (-2,0) -- (-2.8,-0.8) node[below left] {可行 (-)};

% 标记无穷远点
\node at (-1,1) {$\infty$ (平移)};

\node[below] at (-1.5,-0.8) {注意：中间段对应平移，不是旋转};
\end{tikzpicture}
\end{center}

\textbf{为什么中间段"不可行"？}

这里的"不可行"不是指物理上不可能，而是指：
\begin{itemize}
\item 中间段对应$\omega_z \approx 0$的情况
\item 这是\textbf{平移}，不是旋转
\item 平移的CoR在无穷远处，不在有限平面上
\item 所以有限平面上的中间段实际上不对应旋转运动
\end{itemize}

\section{难点3：三个标记的CoR约束}

\subsection{问题设置}

给定静止物体和可动物体之间的接触，接触法线$\hat{n}$指向可动物体内部。

\subsection{标记规则（详细解释）}

\textbf{规则1}：接触法线上的点标记为'$\pm$'

\textbf{为什么？}
\begin{itemize}
\item 如果CoR在接触法线上，接触点处的速度方向与法线垂直
\item 这意味着$F^T V = 0$（满足滚动-滑动约束）
\item 可以是滚动（R）或滑动（S）
\end{itemize}

\textbf{规则2}：法线左侧的点标记为'+'

\textbf{为什么？}
\begin{itemize}
\item 想象你站在接触点，面向法线方向
\item 左侧的点作为CoR时，物体会逆时针旋转
\item 逆时针旋转时，接触点有远离法线的趋势
\item 这满足不可穿透约束：$F^T V \geq 0$
\end{itemize}

\textbf{规则3}：法线右侧的点标记为'-'

\textbf{为什么？}
\begin{itemize}
\item 右侧的点作为CoR时，物体会顺时针旋转
\item 顺时针旋转时，接触点也有远离法线的趋势
\item 同样满足不可穿透约束
\end{itemize}

\subsection{具体例子：单个接触}

\begin{center}
\begin{tikzpicture}[scale=2.5]
% 坐标轴
\draw[->] (-2,0) -- (2,0) node[right] {$x$};
\draw[->] (0,-0.5) -- (0,2.5) node[above] {$y$};

% 接触法线（垂直向上）
\draw[thick,blue,->,line width=2pt] (0,0) -- (0,2) node[above] {$\hat{n}$ (法线)};
\filldraw[black] (0,0) circle (2pt) node[below] {接触点};

% 标记区域
\fill[red!20] (-2,0) -- (-2,2) -- (0,2) -- (0,0) -- cycle;
\node[red] at (-1,1) {\Large +};

\fill[green!20] (0,0) -- (0,2) -- (2,2) -- (2,0) -- cycle;
\node[green] at (1,1) {\Large -};

% 法线上的标记
\draw[thick,blue,line width=2pt] (0,0) -- (0,2);
\node[blue] at (0.3,1) {\Large $\pm$};

% 标记一些具体的CoR
\filldraw[red] (-1,0.5) circle (1.5pt) node[below] {+, B};
\filldraw[red] (-1,1.5) circle (1.5pt) node[above] {+, B};
\filldraw[blue] (0,1) circle (1.5pt) node[right] {$\pm$, R或S};
\filldraw[green] (1,0.5) circle (1.5pt) node[below] {-, B};
\filldraw[green] (1,1.5) circle (1.5pt) node[above] {-, B};

% 添加说明
\node[below] at (0,-0.3) {法线左侧：+，右侧：-，法线上：$\pm$};
\end{tikzpicture}
\end{center}

\subsection{物理解释：为什么左侧是'+'？}

\begin{center}
\begin{tikzpicture}[scale=2]
% 接触点和法线
\filldraw[black] (0,0) circle (2pt) node[below] {接触点};
\draw[thick,blue,->] (0,0) -- (0,1.5) node[above] {$\hat{n}$};

% 左侧的CoR
\filldraw[red] (-1,1) circle (2pt) node[above left] {CoR (+)};

% 从CoR到接触点的向量
\draw[thick,red,dashed] (-1,1) -- (0,0);

% 接触点的速度方向（垂直于CoR到接触点的向量）
\draw[thick,green,->] (0,0) -- (0.3,0.3) node[right] {$v$};

% 标注
\node at (-0.5,0.5) {逆时针};
\node at (0.2,0.2) {速度};
\end{tikzpicture}
\end{center}

当CoR在左侧时：
\begin{itemize}
\item 物体绕CoR逆时针旋转（'+'）
\item 接触点的速度方向向右（垂直于CoR到接触点的方向）
\item 速度与法线的点积：$v \cdot \hat{n} = 0$（垂直）
\item 或者有远离法线的分量，满足$F^T V \geq 0$
\end{itemize}

\section{难点4：多个接触的约束区域}

\subsection{基本思想}

如果有多个接触，每个接触都贡献一个约束区域。可行CoR区域是所有约束区域的\textbf{交集}。

\subsection{例子：两个接触}

\textbf{设置}：
\begin{itemize}
\item 接触1：法线$\hat{n}_1$垂直向上，在$(0, 0)$
\item 接触2：法线$\hat{n}_2$水平向右，在$(2, 1)$
\end{itemize}

\textbf{约束区域}：
\begin{itemize}
\item 接触1：左侧'+'，右侧'-'，法线上'$\pm$'
\item 接触2：下方'+'，上方'-'，法线上'$\pm$'
\end{itemize}

\textbf{可行区域}：两个约束区域的交集

\begin{center}
\begin{tikzpicture}[scale=1.5]
% 接触1的法线
\draw[thick,blue,->] (0,0) -- (0,2);
\filldraw[black] (0,0) circle (2pt) node[below left] {接触1};
\fill[red!10] (-2,0) -- (-2,2) -- (0,2) -- (0,0) -- cycle;
\fill[green!10] (0,0) -- (0,2) -- (2,2) -- (2,0) -- cycle;

% 接触2的法线
\draw[thick,red,->] (2,1) -- (4,1);
\filldraw[black] (2,1) circle (2pt) node[below] {接触2};
\fill[blue!10,rotate around={0:(2,1)}] (2,1) rectangle (4,3);
\fill[yellow!10,rotate around={0:(2,1)}] (2,-1) rectangle (4,1);

% 交集区域（可行区域）
\fill[purple!30] (0,0) -- (0,1) -- (2,1) -- (2,0) -- cycle;
\node[purple] at (1,0.5) {可行区域};

% 标记
\node[red] at (-1,1) {接触1: +};
\node[green] at (1,1) {接触1: -};
\node[blue] at (3,2) {接触2: +};
\node[yellow] at (3,0) {接触2: -};
\end{tikzpicture}
\end{center}

\textbf{分析}：
\begin{itemize}
\item 可行区域必须同时满足两个约束
\item 接触1要求：左侧'+'或右侧'-'
\item 接触2要求：下方'+'或上方'-'
\item 交集：同时满足的区域
\end{itemize}

\subsection{三个接触的例子}

\textbf{关键点}：当有多个接触时，可行区域是\textbf{凸的}。

\textbf{为什么？}
\begin{itemize}
\item 每个接触的约束区域是半空间（凸集）
\item 多个半空间的交集仍然是凸集
\item 因此可行CoR区域是凸的
\end{itemize}

\section{总结：关键理解点}

\subsection{正线性组合的CoR}

\begin{enumerate}
\item \textbf{相同符号}：CoR在两个CoR之间的线段上
\item \textbf{不同符号}：CoR在两条射线上（分别从两个CoR出发），中间段对应平移
\end{enumerate}

\subsection{标记规则}

\begin{enumerate}
\item \textbf{法线左侧}：'+'（逆时针旋转可行）
\item \textbf{法线右侧}：'-'（顺时针旋转可行）
\item \textbf{法线上}：'$\pm$'（滚动或滑动）
\end{enumerate}

\subsection{多个接触}

\begin{enumerate}
\item 每个接触贡献一个约束区域
\item 可行区域是约束区域的交集
\item 可行区域是凸的
\end{enumerate}

\section{超级详细的数值例子}

\subsection{例子：不同符号角速度的完整计算}

让我们用一个具体的数值例子来彻底理解不同符号角速度的情况。

\textbf{设置}：
\begin{align}
V_1 &= (1, 2, 0) \quad \Rightarrow \quad \text{CoR}_1 = (0, 2), \quad \text{标记：'+'} \\
V_2 &= (-1, 0, 2) \quad \Rightarrow \quad \text{CoR}_2 = (2, 0), \quad \text{标记：'-'}
\end{align}

\textbf{正线性组合}：$V = k_1 V_1 + k_2 V_2$，其中$k_1, k_2 \geq 0$

计算：
\begin{align}
V &= k_1(1, 2, 0) + k_2(-1, 0, 2) \\
&= (k_1 - k_2, 2k_1, 2k_2)
\end{align}

\textbf{角速度}：$\omega_z = k_1 - k_2$

\textbf{CoR位置}（当$\omega_z \neq 0$时）：
\begin{equation}
\text{CoR} = \left(-\frac{2k_2}{k_1 - k_2}, \frac{2k_1}{k_1 - k_2}\right)
\end{equation}

\textbf{情况分析}：

\begin{enumerate}
\item \textbf{当$k_1 = 2, k_2 = 0$时}：
  \begin{itemize}
  \item $\omega_z = 2 - 0 = 2 > 0$ → 标记：'+'
  \item $\text{CoR} = (0, 2) = \text{CoR}_1$
  \item 这就是$V_1$本身
  \end{itemize}

\item \textbf{当$k_1 = 1.5, k_2 = 0.5$时}：
  \begin{itemize}
  \item $\omega_z = 1.5 - 0.5 = 1 > 0$ → 标记：'+'
  \item $\text{CoR} = \left(-\frac{1}{1}, \frac{3}{1}\right) = (-1, 3)$
  \item CoR在从CoR$_1$出发的射线上
  \end{itemize}

\item \textbf{当$k_1 = 1, k_2 = 1$时}：
  \begin{itemize}
  \item $\omega_z = 1 - 1 = 0$ → 纯平移
  \item CoR在无穷远处
  \item 这是两个CoR之间的"分界点"
  \end{itemize}

\item \textbf{当$k_1 = 0.5, k_2 = 1.5$时}：
  \begin{itemize}
  \item $\omega_z = 0.5 - 1.5 = -1 < 0$ → 标记：'-'
  \item $\text{CoR} = \left(-\frac{3}{-1}, \frac{1}{-1}\right) = (3, -1)$
  \item CoR在从CoR$_2$出发的射线上
  \end{itemize}

\item \textbf{当$k_1 = 0, k_2 = 2$时}：
  \begin{itemize}
  \item $\omega_z = 0 - 2 = -2 < 0$ → 标记：'-'
  \item $\text{CoR} = (2, 0) = \text{CoR}_2$
  \item 这就是$V_2$本身
  \end{itemize}
\end{enumerate}

\textbf{关键观察}：

\begin{itemize}
\item 当$k_1 > k_2$时：$\omega_z > 0$，CoR标记为'+'，在从CoR$_1$出发的射线上
\item 当$k_1 = k_2$时：$\omega_z = 0$，纯平移，CoR在无穷远
\item 当$k_1 < k_2$时：$\omega_z < 0$，CoR标记为'-'，在从CoR$_2$出发的射线上
\end{itemize}

\textbf{图示}：

\begin{center}
\begin{tikzpicture}[scale=1.2]
% 两个CoR
\filldraw[red] (0,2) circle (3pt) node[above left] {CoR$_1$ (+)};
\filldraw[blue] (2,0) circle (3pt) node[below right] {CoR$_2$ (-)};

% 连接直线
\draw[thick,green] (-2,4) -- (4,-2);

% 标记不同区域
\draw[thick,red,->,line width=2pt] (0,2) -- (-1.5,3.5) node[above left] {$k_1 > k_2$\\$\omega_z > 0$\\标记：'+'};

\draw[thick,blue,->,line width=2pt] (2,0) -- (3.5,-1.5) node[below right] {$k_1 < k_2$\\$\omega_z < 0$\\标记：'-'};

% 标记中间点（平移）
\node[purple] at (1,1) {$k_1 = k_2$\\$\omega_z = 0$\\平移（$\infty$）};

% 标记一些具体点
\filldraw[red] (-1,3) circle (2pt) node[above] {$(-1, 3)$};
\filldraw[blue] (3,-1) circle (2pt) node[below] {$(3, -1)$};

\node[below] at (1,-1) {注意：中间段对应平移，不在有限平面上};
\end{tikzpicture}
\end{center}

\section{常见误解澄清}

\subsection{误解1：中间段"不可行"是什么意思？}

\textbf{澄清}：
\begin{itemize}
\item 不是物理上不可能
\item 而是指中间段对应$\omega_z \approx 0$（平移）
\item 平移的CoR在无穷远处，不在有限平面上
\item 所以有限平面上的中间段不对应旋转运动
\item 从数学上看，当$\omega_z \to 0$时，CoR位置$\to \infty$
\end{itemize}

\subsection{误解2：为什么不同符号时是射线而不是线段？}

\textbf{澄清}：
\begin{itemize}
\item 因为$\omega_z$的符号会改变
\item 当$k_1$和$k_2$变化时，$\omega_z$从正变负（或从负变正）
\item 中间有一个点对应$\omega_z = 0$（平移）
\item 所以可行区域被分成两部分：一条'+'射线和一条'-'射线
\end{itemize}

\subsection{误解3：标记'+'和'-'的含义}

\textbf{澄清}：
\begin{itemize}
\item '+'不是"正方向"，而是"逆时针旋转"
\item '-'不是"负方向"，而是"顺时针旋转"
\item 标记表示旋转方向，不是位置的正负
\end{itemize}

\end{document}

